\subsection{7.4  Key Findings}

The simulation produces four findings that directly validate the APIP framework's design decisions.

Finding 1: Tier mismatch is the dominant source of excess cost in ad-hoc remediation. Across all scenarios, ad-hoc responders applied Tier 1 (prompt-level) interventions as a first response regardless of failure mode. In the Behavioral Drift and Alignment Creep scenarios, this tier mismatch produced a secondary regression event (captured in the regression rate column of Table 3) that extended total time-to-resolution by a factor of 2.1$\times$ on average. The APIP's cause attribution phase, by mandating failure mode classification before intervention selection, eliminates this error class entirely in the simulation.

Finding 2: Input Brittleness is invisible to aggregate monitoring. The Input Brittleness scenario is injected at tick 10 but not detected until tick 24 under aggregate monitoring. The billing dispute category's task completion rate drops to 0.41 at injection---well below the 0.82 threshold---while the overall rate remains at 0.83, masking the failure. The APIP's requirement for disaggregated cause attribution analysis detects the failure at tick 10 in the structured condition. This 14-tick detection delay translates to a $-$0.61 CSAT impact on the affected user segment, representing a meaningful customer experience gap invisible to traditional monitoring.

Finding 3: Token consumption and latency are leading indicators of degradation. In all four failure modes, token consumption and p95 latency begin increasing before any quality metric crosses its absolute threshold. In the Behavioral Drift scenario, token consumption increases by 18\% and latency increases by 120ms before the policy citation error rate trigger fires. Rath's [24] Agent Stability Index independently identifies tool usage patterns and reasoning pathway stability as early degradation signals. Including token and latency drift as soft APIP triggers could reduce the injection-to-detection gap and initiate cause attribution earlier.

Finding 4: CSAT impact is failure-mode-specific and differentially lagged. Tool Misuse produces the smallest peak CSAT impact ($-$0.19) because customer dissatisfaction is triggered by incorrect actions rather than incorrect responses, and many tool errors are resolved through retry without visible customer impact. Alignment Creep produces a gradual CSAT decline ($-$0.31 at detection) because tone and policy drift affect satisfaction cumulatively. Input Brittleness produces the largest segment-specific impact ($-$0.61) because affected customers encounter categorical failure rather than degraded performance. These differential profiles validate the taxonomy's distinction between failure modes and support mode-specific monitoring sensitivity settings.

